\Needspace{10\baselineskip}
\item \textbf{Ordenamiento y busqueda en una lista de periodos orbitales}.
\nopagebreak[4]

\begin{tcolorbox}[colback=blue!2!white,colframe=blue!55!black,title={Enunciado},
                  before skip=3pt,after skip=5pt,boxsep=3pt,top=3pt,bottom=3pt,
                  breakable]
\small
\raggedright
En una observacion simplificada se registraron los periodos orbitales, en dias, de
un conjunto de exoplanetas:
\[
\texttt{periodos = [12.4,\ 3.1,\ 8.7,\ 25.0,\ 1.9,\ 14.2,\ 6.5,\ 18.3]}.
\]

Utilice los siguientes algoritmos ya dados. No debe construirlos desde cero:

\begin{pythonjupyter}
def bubble_sort(lista):
    n = len(lista)
    for i in range(n - 1):
        for j in range(n - 1 - i):
            if lista[j] > lista[j + 1]:
                lista[j], lista[j + 1] = lista[j + 1], lista[j]
    return lista

def binary_search(lista_ordenada, objetivo):
    low, high = 0, len(lista_ordenada) - 1

    while low <= high:
        mid = (low + high) // 2

        if lista_ordenada[mid] == objetivo:
            return mid
        elif objetivo < lista_ordenada[mid]:
            high = mid - 1
        else:
            low = mid + 1

    return -1
\end{pythonjupyter}

Realice las siguientes operaciones:

\begin{itemize}
\setlength{\topsep}{4pt}
\setlength{\itemsep}{3pt}
\setlength{\parsep}{0pt}
\item[(a)] Defina la lista \texttt{periodos}.
\item[(b)] Ordene los periodos de menor a mayor usando \texttt{bubble\_sort}. Use una copia de la lista original para no modificarla directamente.
\item[(c)] Imprima la lista original y la lista ordenada.
\item[(d)] Determine el periodo minimo y el periodo maximo usando la lista ordenada.
\item[(e)] Use \texttt{binary\_search} para buscar el periodo \texttt{14.2} en la lista ordenada e imprima su posicion.
\item[(f)] Use \texttt{binary\_search} para buscar el periodo \texttt{10.0}. Indique con un mensaje si no fue encontrado.
\item[(g)] Explique brevemente, dentro de un \texttt{print}, por que la busqueda binaria debe aplicarse sobre una lista ordenada.
\end{itemize}
\end{tcolorbox}

\begin{solution}
\begin{pythonjupyter}
# (a)
periodos = [12.4, 3.1, 8.7, 25.0, 1.9, 14.2, 6.5, 18.3]

# Algoritmos entregados en el enunciado
def bubble_sort(lista):
    n = len(lista)
    for i in range(n - 1):
        for j in range(n - 1 - i):
            if lista[j] > lista[j + 1]:
                lista[j], lista[j + 1] = lista[j + 1], lista[j]
    return lista

def binary_search(lista_ordenada, objetivo):
    low, high = 0, len(lista_ordenada) - 1

    while low <= high:
        mid = (low + high) // 2

        if lista_ordenada[mid] == objetivo:
            return mid
        elif objetivo < lista_ordenada[mid]:
            high = mid - 1
        else:
            low = mid + 1

    return -1

# (b)
periodos_ordenados = bubble_sort(periodos.copy())

# (c)
print("Lista original:", periodos)
print("Lista ordenada:", periodos_ordenados)

# (d)
periodo_minimo = periodos_ordenados[0]
periodo_maximo = periodos_ordenados[-1]

print("Periodo minimo:", periodo_minimo)
print("Periodo maximo:", periodo_maximo)

# (e)
posicion_142 = binary_search(periodos_ordenados, 14.2)
print("Posicion de 14.2:", posicion_142)

# (f)
posicion_100 = binary_search(periodos_ordenados, 10.0)

if posicion_100 == -1:
    print("El periodo 10.0 no fue encontrado")
else:
    print("Posicion de 10.0:", posicion_100)

# (g)
print("La busqueda binaria necesita una lista ordenada porque en cada paso")
print("compara el objetivo con el elemento central y descarta una mitad.")
print("Si los datos no estan ordenados, esa decision no es valida.")
\end{pythonjupyter}
\end{solution}
